\section{General approach}

We will examine three different variations of mesh intersections. To illustrate these situations, we will consider a finite element model of a human head. Such models contain a triangular surface mesh for each tissue boundary, e.g., skull and brain, based on the segmentation of an MRI scan. The human brain presents a modeling challenge since its outer surface consists of ridges called gyri that are closer together than the width of a single MRI voxel. If two neighboring gyri intersect, we can choose to either merge or detach the two gyri, depending on the application. A different situation that causes intersections is when two enclosed meshes intersect with each other, for example, in areas where the brain and skull touch. In this study, we represent these cases with simple geometries that are easier to visualize and quantify. For example, we used two individual spheres that intersect on one side to represent two gyri on one surface intersecting (Fig. \ref{fig:teaser}, first image), and two enclosed spheres to represent a brain and skull surface intersecting (Fig. \ref{fig:sol2-2}, left). These simple geometries closely represent the local problems encountered in large complex structures. For example, a brain surface may have many self-intersecting gyri, each of which resembles the case of two spheres, and can be repaired in the same way. To ensure a comprehensive comparison, test meshes of varying complexity were generated using Blender. We created three test cases with two neighboring meshes and one with two enclosed meshes.

In all cases, we aimed to achieve three goals: 1) fix all intersections without producing other defects such as holes, 2) produce a high-quality mesh, and 3) preserve the original shape and position of the mesh as much as possible. To accomplish this, we can either locally resolve intersecting areas, or we can globally remesh the model. We applied and compared five different methods to merge two neighboring meshes. This approach is sufficient for most cases and is commonly adopted in existing repair functions across various mesh libraries. We also developed one method to detach two neighboring meshes and one method to detach two enclosed meshes. Ideally, the final output in each case should have no intersections and no significant shape alterations. 

To quantitatively assess the output meshes, we  evaluated them using the mesh metrics and statistics shown in Table 1 in the supplemental material. In this study, we place greater emphasis on mesh quality than on computational efficiency, as we evaluate the final mesh condition. Ideally, a high-quality triangular mesh consists of elements that are as close to equilateral as possible. Each metric was evaluated based on this assumption. Metrics marked with an asterisk were computed using PyVista, a wrapper of VTK. The rest were calculated manually. Each test mesh contains noise, in the form of self-intersections. These evaluations were performed on a computer with an 11th Gen Intel Core CPU and 12.0 GB of RAM. 
